\documentclass[UTF8,a4paper,11pt]{ctexart}
\usepackage[margin=27mm]{geometry}
\usepackage{amsmath,amssymb,amsthm,mathtools,booktabs,hyperref,array}
\hypersetup{colorlinks=true,linkcolor=blue,citecolor=blue,urlcolor=blue}
\numberwithin{equation}{section}
\newtheorem{theorem}{定理}[section]
\newtheorem{lemma}[theorem]{引理}
\newtheorem{proposition}[theorem]{命题}
\theoremstyle{definition}
\newtheorem{definition}[theorem]{定义}
\theoremstyle{remark}
\newtheorem{remark}[theorem]{说明}
\newcommand{\Ftwo}{\mathbb F_2}
\newcommand{\tr}{\operatorname{Tr}}
\newcommand{\dimf}{\dim_{\Ftwo}}
\newcommand{\Psub}{P_E}
\title{卡莱特零和子空间猜想的证明及其 Lean 形式化}
\author{王泉智\\[-.3em]\small\href{mailto:mrhwrn120504@gmail.com}{\texttt{mrhwrn120504@gmail.com}}\and ChatGPT\\[-.3em]\small AI 研究协作者}
\date{2026年11月}

\begin{document}
\maketitle

\begin{abstract}
设 $K$ 是含 $2^n$ 个元素的有限域，并把它看作 $\Ftwo$-向量空间。对 $x\ne0$，考虑乘法逆映射 $x\mapsto x^{-1}$。本文证明 Carlet 零和子空间猜想的奇偶次数完整范围：若 $n$ 为奇数且 $3\le k,\ k+3\le n$，或 $n$ 为偶数且 $2\le k,\ k+2\le n$，则存在 $k$ 维 $\Ftwo$-线性子空间 $E\subseteq K$，使
\[
  \sum_{x\in E\setminus\{0\}}x^{-1}=0.
\]
奇数次数的证明以根积多项式、椭圆曲线 Hasse 界、低维构造和高维能量迭代为基础；偶数次数的证明由二、三维基例、维数递推和互补维数对称性给出。两部分均已在 Lean 4 中对相应量词及任意有限域完整形式化。完整默认构建、独立验证和主定理公理审计均已通过：无编译警告，源码与 OLean 哈希一致，依赖只含 \texttt{propext}、\texttt{Classical.choice} 和 \texttt{Quot.sound}，不含 \texttt{sorryAx} 或自建公理。
\end{abstract}

\noindent\textbf{关键词：} 有限域；乘法逆映射；零和子空间；子空间多项式；Kloosterman 和；Hasse 界；Lean 形式化

\section{引言与主定理}

\noindent\textbf{与此前研究的关系。} Cheng 于 2026 年 10 月的预印本~\cite{Cheng2026} 已证明 Carlet 的存在性猜想，并给出定量计数定理。本文的贡献在于 Lean 4 形式化、显式构造及可审计的机器验证；我们不主张对原存在性猜想的首次证明优先权。


设 $K$ 是特征为 $2$ 的有限域。其非零元素构成有限乘法群，但乘法逆映射 $x\mapsto x^{-1}$ 并非加法映射。因此，一个自然的有限域问题是：给定维数时，是否必然存在一个加法子空间，使其中所有非零元素的逆元之和恰好抵消为零？本文研究 Carlet 工作中的这一零和问题\cite{Carlet2025}，并对固定基数的所有有限域统一给出结论。

对 $\Ftwo$-线性子空间 $E\subseteq K$，记
\[
 S(E)=\sum_{x\in E\setminus\{0\}}x^{-1}.
\]
本文先分别给出奇数与偶数次数的适用维数范围，再将它们拼接为完整的奇偶次数结论。Lean 形式化对两部分分别保留原始量词。

\begin{definition}
若 $E$ 是 $K$ 的有限维 $\Ftwo$-线性子空间，称 $E$ 为\emph{零和子空间}，若
\[
  S(E):=\sum_{x\in E\setminus\{0\}}x^{-1}=0.
\]
\end{definition}

\begin{theorem}[奇数次数 Carlet 命题]\label{thm:main}
设 $n,k\in\mathbb N$，$n$ 为奇数，$3\le k$ 且 $k+3\le n$。对任意特征为 $2$ 且基数为 $2^n$ 的有限域 $K$，存在 $k$ 维 $\Ftwo$-线性子空间 $E\subseteq K$，使 $S(E)=0$。
\end{theorem}

本文证明对有限域 $K$ 的选择是统一的：Lean 定理对所有满足字段结构、有限性、$\Ftwo$-代数结构及基数等式的 $K$ 同时量化。奇数部分的机器陈述见第\,\ref{sec:formal}节；偶数部分随后单独陈述。两个适用范围各自保留其基例条件，不作渐近外推。

\section{根积多项式与系数判据}

对有限 $\Ftwo$-子空间 $E\subseteq K$，定义
\[
  P_E(T)=\prod_{u\in E}(T+u).
\]
由于 $E$ 对加法封闭，$P_E$ 是线性化多项式。若 $r=\dimf E$，则存在系数 $c_i\in K$ 使
\[
  P_E(T)=c_0T+c_1T^2+c_2T^4+\cdots+T^{2^r},\qquad c_0\ne0.
\]

\begin{lemma}[倒数和判据]\label{lem:coeff}
对任意非零有限维 $\Ftwo$-子空间 $E$，有
\[
  S(E)=\frac{c_1}{c_0},\qquad\text{因而}\qquad S(E)=0\Longleftrightarrow [T^2]P_E=0.
\]
\end{lemma}
\begin{proof}
将 $P_E(T)$ 中的零根因子分离：
\[
 P_E(T)=T\prod_{u\in E\setminus\{0\}}(T+u).
\]
右侧乘积的常数项为 $\prod_{u\ne0}u=c_0$，一次项为
\[
 \sum_{u\ne0}\prod_{v\ne0,u}v
 =\Bigl(\prod_{v\ne0}v\Bigr)\sum_{u\ne0}u^{-1}=c_0S(E).
\]
乘以 $T$ 后，这正是 $P_E$ 的 $T^2$ 系数 $c_1$。由于 $c_0\ne0$，结论成立。
\end{proof}

设 $I(E)=P_E(K)$，并令 $D(E)=I(E)^\perp$，其中正交补相对于迹配对。此时 $\dimf I(E)=n-\dimf E$，而 $\dimf D(E)=\dimf E$。若 $\dimf E=k$，记 $a_{k-1}=[T^{2^{k-1}}]P_E$，则形式化的对偶多项式恒等式给出
\begin{equation}\label{eq:duality}
 S(D(E))=a_{k-1}^{2}.
\end{equation}
所以寻找一个 $k$ 维空间使其根积多项式的次高 Frobenius 系数为零，等价于寻找一个 $k$ 维零和子空间。这是后文从系数消失回到倒数和的桥梁；通常该零和空间是 $D(E)$，不必等于 $E$。

若 $x\notin E$，记 $E+\langle x\rangle$ 为扩张一维的子空间。根按两个陪集 $E$ 与 $x+E$ 分组，利用 $P_E(T+x)=P_E(T)+P_E(x)$ 得到
\begin{equation}\label{eq:extension}
 P_{E+\langle x\rangle}(T)=P_E(T)^2+P_E(x)P_E(T).
\end{equation}
式\eqref{eq:extension}是后面所有小维构造的基本递推。

\begin{lemma}[互补维数对称性]\label{lem:complement}
若 $K$ 的 $\Ftwo$-维数为 $n$ 且 $|K|>2$，则对 $0\le k\le n$，$K$ 中存在 $k$ 维零和子空间，当且仅当存在 $(n-k)$ 维零和子空间。
\end{lemma}
\begin{proof}
映射 $P_E:K\to K$ 是 $\Ftwo$-线性的，核为 $E$，像 $I(E)$ 的维数为 $n-k$。子空间多项式满足
\[
 P_{I(E)}(P_E(T))=T^q-T.
\]
写 $c_0=[T]P_E\ne0$、$c_1=[T^2]P_E$、$d_0=[T]P_{I(E)}\ne0$、$d_1=[T^2]P_{I(E)}$。比较两侧 $T^2$ 系数得到 $d_0c_1+d_1c_0^2=0$，故 $c_1=0$ 当且仅当 $d_1=0$。结合引理\,\ref{lem:coeff}，$E$ 零和当且仅当 $I(E)$ 零和。又因为 $\dimf I(E)=n-k$，便得结论；反向由同一映射的互补维数等式应用一次即可。
\end{proof}

\section{三维基例：椭圆曲线与 Kloosterman 和}

令 $q=|K|=2^n$，并定义二元 Kloosterman 和
\[
 \operatorname{Kl}_K=\sum_{v\in K^\times}(-1)^{\tr(v+v^{-1})}.
\]
考虑特征为 $2$ 的 Weierstrass 曲线
\[
 C:\quad y^2+xy=x^3+1.
\]
其判别式为 $1$，故是椭圆曲线。按 $x$ 对仿射点分纤维计数，并利用 Artin--Schreier 方程 $z^2+z=c$ 有 $1+(-1)^{\tr(c)}$ 个解，得到
\[
 \#C(K)=q+1+\operatorname{Kl}_K.
\]
对该椭圆曲线应用 Hasse 界\cite{Silverman} $|\#C(K)-(q+1)|\le2\sqrt q$，便有 $|\operatorname{Kl}_K|\le2\sqrt q$。

对所有含 $1$ 的三维子空间按其下一 Frobenius 系数 $b$ 分组。Lean 中的精确计数恒等式为
\begin{equation}\label{eq:p3}
 (q-2)(q-4)\,p_3(0)=q+1-3\operatorname{Kl}_K,
\end{equation}
其中 $p_3(0)$ 是相应有限集合中系数为零的比例。若 $n$ 是奇数且 $n\ge7$，则 $q\ge128$，从而
\[
 q+1-3\operatorname{Kl}_K\ge q+1-6\sqrt q>0.
\]
式\eqref{eq:p3}左边的因子 $(q-2)(q-4)$ 为正，故 $p_3(0)>0$。因此存在三维子空间使 $[T^2]P_E=0$，由引理\,\ref{lem:coeff}得到三维零和子空间。

\begin{proposition}\label{prop:k3}
任意奇数 $n\ge7$ 及任意 $|K|=2^n$ 的有限域 $K$ 都含有三维零和子空间。
\end{proposition}

\section{四维基例：范数一曲线上的 Fourier 计数}

本节给出覆盖所有适用奇数次数的四维构造。大次数证明的关键不是单独寻找一个曲线点，而是同时施加两个迹正交条件，确保后续两个值都落入指定的子空间多项式像。

\subsection{二维起点和扩维恒等式}

取 $t\in K\setminus\{0,1\}$，令 $H=\langle1,t\rangle_{\Ftwo}$。因 $n$ 为奇数，$T^2+T+1$ 在 $K$ 中无根，故 $c=t^2+t+1\ne0$。直接展开得
\[
 P_H(T)=T^4+cT^2+(t^2+t)T.
\]
考虑范数一曲线
\begin{equation}\label{eq:normcurve}
 \mathcal C(K)=\{(a,u)\in K^2:a^2+au+u^2=1\}.
\end{equation}
当 $n$ 为奇数时，$\#\mathcal C(K)=q+1$。构造一个二维迹平面 $D$，使 $a,u\in D^\perp$ 可分别写成
\[
 P_H(x)=c^2a,\qquad P_H(y)=c^2u.
\]
具体地，令 $L(x)=c^{-2}P_H(x)$，取 $D=(\operatorname{im}L)^\perp$，其中正交补相对于非退化迹配对 $(r,s)\mapsto\tr(rs)$。因为 $\ker L=H$，所以 $\dimf\operatorname{im}L=n-2$、$\dimf D=2$，且 $D^\perp=\operatorname{im}L$。这保证了所需的 $x,y$ 存在。

对曲线点要求 $u\in D^\perp$ 且 $a+u\in D^\perp$，即
\[
 \tr(ru)=0,\quad\tr(s(a+u))=0\qquad(r,s\in D).
\]
记 $\chi(z)=(-1)^{\tr(z)}$。对频率 $(A,B)\ne(0,0)$，Fourier 系数可化为
\[
 \mathcal A(A,B)=\sum_{x\in K}\chi\!\left(\frac{Bx^2+A}{x^2+x+1}\right)+\chi(B).
\]
其计数来源是 Artin--Schreier 方程
\[
 \delta^2+\delta=\frac{(A+B)x^2+Ax}{x^2+x+1}.
\]
把一般解通过双有理坐标变换送到 Weierstrass 曲线
\[
 y^2+xy+By=x^3+(A^2+B+1)x^2+(A+(A^2+B))x+B(A+(A^2+B))+A^2.
\]
该曲线判别式为 $(A^2+AB+B^2)^2$。奇数次数保证 $T^2+T+1$ 在 $K$ 中无根，也保证非零频率下此判别式非零。Hasse 界控制曲线总点数；一般点与仿射方程解之间的映射是单射，两个退化纤维总计至多损失五点。因此得到每个非零频率的下界 $\mathcal A(A,B)\ge-2\sqrt q-4$。用 $D$ 的四个元素分别对两个条件作加法特征过滤，得到 16 项 Fourier 展开。零频项为 $q+1$；其余 15 项满足上述下界。若 $N$ 表示满足两个迹条件的曲线点数，便有
\begin{equation}\label{eq:k4count}
 16N\ge q+1-30\sqrt q-60.
\end{equation}
当 $q\ge2048$ 时，右侧严格大于 $48$，故 $N>3$。

由 $u,a+u\in D^\perp$ 及子空间封闭性可知 $a\in D^\perp$。范数方程下，若 $a=0$，则 $u=1$；若 $a=1$，则 $u=0$ 或 $u=1$。所以退化点至多三个。由 $N>3$ 可取满足迹条件且 $a\ne0,1$ 的 $(a,u)$。迹正交性给出前述 $x,y$。

令 $V=H+\langle x\rangle$。由\eqref{eq:extension}，$\dimf V=3$，并且
\[
 [T^4]P_V=c^2+P_H(x)=c^2(1+a)\ne0,
\]
而
\[
 P_V(y)=P_H(y)^2+P_H(x)P_H(y)
       =c^4(u^2+au)=c^4(1+a^2)=\bigl(c^2(1+a)\bigr)^2.
\]
最后一步使用 $a^2+au+u^2=1$ 和特征 $2$ 下 $(1+a)^2=1+a^2$。因 $a\ne1$，上述 $T^4$ 系数非零；因此 $y\notin V$。将 $V$ 再扩张 $y$ 后，式\eqref{eq:extension}说明四维根积的 $T^8$ 系数为零。通过引理\,\ref{lem:coeff}的同维系数—倒数和接口，得到四维零和子空间。

\begin{proposition}\label{prop:k4large}
若 $n$ 为奇数且 $n\ge11$，则任意 $|K|=2^n$ 的 $K$ 含有四维零和子空间。
\end{proposition}

当 $n=7,9,11,13,15$ 时，形式化开发另给出精确有限域构造；小次数分支与上述大次数分支拼接后得到：
\begin{theorem}\label{thm:k4}
对每个奇数 $n\ge7$ 和任意 $|K|=2^n$ 的有限域 $K$，均存在四维零和子空间。
\end{theorem}

\section{五维基例：从三维种子到范数曲线}

\subsection{由四维结果制造非零系数种子}

由定理\,\ref{thm:k4}及互补系数判据，取四维子空间 $W$ 使 $[T^8]P_W=0$。取三维超平面 $U<W$，并取 $v\in W\setminus U$。令 $b=[T^4]P_U$。由\eqref{eq:extension}比较 $T^8$ 系数可得
\[
 0=[T^8]P_W=b^2+P_U(v).
\]
根 $P_U$ 恰为 $U$，而 $v\notin U$，所以 $P_U(v)\ne0$，从而 $b\ne0$。置 $c=b^2\ne0$。

令 $L=c^{-1}P_U$，其像的余维为三。取 $D=(\operatorname{im}L)^\perp$，于是 $\dimf D=3$ 且 $|D|=8$。对范数一曲线\eqref{eq:normcurve}上的点施加 $u\in D^\perp$ 与 $a+u\in D^\perp$ 两个条件。对 $D$ 的八个元素分别过滤两项，得到 64 项 Fourier 展开；零频项为 $q+1$，其他 63 项每项下界为 $-2\sqrt q-4$。因此合格点数 $N$ 满足
\begin{equation}\label{eq:k5count}
 64N\ge q+1-63(2\sqrt q+4)=q+1-126\sqrt q-252.
\end{equation}
若 $q\ge65536$，右侧严格大于 $192$，故 $N>3$。曲线上 $a=0$ 至多给出一个点，而 $a=1$ 至多给出两个点，故可取合格点 $(a,u)$ 满足 $a\ne0,1$。

由 $u,a+u\in D^\perp=\operatorname{im}L$，相加得 $a\in D^\perp$。故存在 $x,y\in K$ 使
\[
 P_U(x)=ca,\qquad P_U(y)=cu.
\]
因为 $a,c\ne0$，有 $x\notin U$。令 $V=U+\langle x\rangle$，于是 $\dimf V=4$。记 $A=[T^8]P_V$。由\eqref{eq:extension}以及 $P_U$ 首一、次数为 $8$，可得
\[
 A=b^2+P_U(x)=c(1+a)\ne0.
\]
又由范数方程 $u^2+au=1+a^2$，
\[
 P_V(y)=P_U(y)^2+P_U(x)P_U(y)
       =c^2(u^2+au)=c^2(1+a^2)=A^2\ne0.
\]
因此 $y\notin V$。扩张成五维空间 $Z=V+\langle y\rangle$ 后，式\eqref{eq:extension}给出
\[
 [T^{16}]P_Z=A^2+P_V(y)=0.
\]
子空间多项式的同维对偶系数判据将这一零系数等价地转化为五维零和子空间。这里不能省略这一步：五维的 $T^{16}$ 系数判据并非直接等同于倒数和判据，而是经由已证明的系数对偶定理相连。

\begin{proposition}\label{prop:k5large}
若 $n$ 为奇数且 $n\ge17$，则任意 $|K|=2^n$ 的 $K$ 含有五维零和子空间。
\end{proposition}

剩余允许的小次数为 $n=9,11,13,15$。项目用精确的有限域模型和可检查的线性代数证书构造相应五维空间，并通过有限域同构将证书转移到任意基数 $2^n$ 的有限域。故有：
\begin{theorem}\label{thm:k5}
对每个奇数 $n\ge9$ 和任意 $|K|=2^n$ 的有限域 $K$，均存在五维零和子空间。
\end{theorem}

\section{高维：实际概率递推与能量收缩}

令 $\Omega_d$ 表示所有包含 $1$ 的 $d$ 维 $\Ftwo$-子空间。对 $E\in\Omega_d$，记 $b(E)$ 为 $P_E$ 的下一 Frobenius 系数；定义其精确概率分布
\[
 p_d(b)=\frac{\#\{E\in\Omega_d:b(E)=b\}}{\#\Omega_d},\qquad
 f_d(b)=p_d(b)-q^{-1}.
\]
于是 $\sum_{b\in K}f_d(b)=0$。形式化证明从实际有限集合计数建立子空间从 $d$ 维到 $d+1$ 维的递推，而不是把一个抽象随机模型当成证明。

给定非平凡加法特征 $\psi:K\to\mathbb R$ 和实数 $s$，对零和函数 $f:K\to\mathbb R$ 定义转移算子
\[
 (T_sf)(b)=\frac{1}{q-s}\begin{cases}
 (q-s)f(0),&b=1,\\[2pt]
 \displaystyle\sum_{v\ne0}\bigl(1+(s-1)\psi((b-1)/v)\bigr)f(v),&b\ne1.
 \end{cases}
\]
该定义使用平移坐标 $b-1$，特殊分支 $b=1$ 对应零坐标。Lean 源码 \texttt{TransferOperator.lean} 给出该算子的正式定义；实际概率递推已证明与它一致。令能量 $\mathcal E(f)=\sum_b f(b)^2$。

\begin{lemma}[两步能量收缩]\label{lem:energy}
若 $s\ge2$ 且 $32s^2\le q$，则对零和函数 $f$，
\[
 \mathcal E(T_{2s}(T_sf))\le\frac{4s^2}{q}\mathcal E(f).
\]
当 $s^2\le q$ 且 $s<q$ 时，一步转移满足 $\mathcal E(T_sf)\le\mathcal E(f)$。
\end{lemma}
\begin{proof}
把 $T_s$ 拆成关于特征值方向的秩一修正。有限加法特征的正交关系给出精确平方范数公式；对两步复合，主方向系数可界为 $2s/\sqrt q$，其平方就是 $4s^2/q$。假设 $32s^2\le q$ 保证分母 $q-s,q-2s$ 为正，并控制剩余秩一项；同一公式在 $s^2\le q$ 下给出一步非扩张。Lean 中先证明带全部异常方向项的精确恒等式，再作上述不等式估计。
\end{proof}

初始的一维分布是点质量，其中心化能量为 $1-q^{-1}$。将引理\,\ref{lem:energy}沿奇数维度逐次应用，得到对于 $k=2r+1$：
\[
 \mathcal E(f_k)\le \frac{2^{2r^2+2r}}{q^r}(1-q^{-1}).
\]
对 $k=2r+2$，先用该奇数维估计，再用一次能量非扩张，得到同一上界。若 $p_k(0)=0$，则 $f_k(0)=-q^{-1}$，从而单个坐标已贡献 $q^{-2}$，即 $\mathcal E(f_k)\ge q^{-2}$。故只要上界严格小于 $q^{-2}$，必有 $p_k(0)>0$。

指数比较由条件给出。奇数维情形 $r\ge5$ 且 $n\ge4r+3$ 时，
\[
 2r^2+2r<n(r-2),
\]
所以 $2^{2r^2+2r}/q^r<q^{-2}$。偶数维情形 $r\ge4$ 且 $n\ge4r+5$ 时同样有该严格不等式。因此得到 $k\ge11$ 奇数或 $k\ge10$ 偶数、且 $n\ge2k+1$ 时的高维存在性。$p_k(0)>0$ 意味着某个实际子空间的下一系数为零；经同维对偶和引理\,\ref{lem:coeff}即为所求零和子空间。

\section{偶数次数：二维基例、递推与互补维数}

偶数次数部分的起点与奇数部分不同。若 $n$ 为偶数，则 $K$ 含有唯一的四元子域 $\mathbb F_4$。取 $\omega\in\mathbb F_4\setminus\mathbb F_2$，有 $\omega^2+\omega+1=0$；该子域的三个非零元素的逆元之和为 $1+\omega+\omega^2=0$，因此得到 $k=2$ 基例。对偶数 $n\ge6$，三维基例由 Kloosterman 和与椭圆曲线 Hasse 界得到。

\begin{theorem}[偶数次数命题]\label{thm:even}
设 $n$ 为偶数，$|K|=2^n$。若 $2\le k$ 且 $k+2\le n$，则 $K$ 含有 $k$ 维零和子空间。
\end{theorem}
\begin{proof}
先说明三维基例。令 $q=|K|$，并记 $\operatorname{Kl}_K=\sum_{v\in K^\times}(-1)^{\operatorname{Tr}(v+v^{-1})}$。椭圆曲线 $y^2+xy=x^3+1$ 的点数给出 $|\operatorname{Kl}_K|\le2\sqrt q$。对包含 $1$ 的三维子空间按下一 Frobenius 系数分组，偶数次数的精确计数恒等式为
\[
 (q-2)(q-4)p_3(0)=q-13+3\operatorname{Kl}_K.
\]
这里偶数次数下 $\operatorname{Tr}(1)=0$，所以对应的加法特征加权计数与奇数次数公式不同。若 $n\ge6$，则 $q\ge64$，从 Hasse 界得到
\[
q-13+3\operatorname{Kl}_K\ge q-13-6\sqrt q=(\sqrt q-3)^2-22\ge3>0.
\]
又 $(q-2)(q-4)>0$，故 $p_3(0)>0$，确有三维系数零的子空间；系数判据及同维对偶给出三维零和子空间。

再考虑低半维数 $2k\le n$。偶数维数从二维基例开始，每次由 $k-2$ 维零和空间与二维零和空间作维数相加；奇数维数从上述三维基例开始，每次由 $k-2$ 维和二维空间相加。相加引理要求两维数之积小于 $n$，此处为 $2(k-2)<n$，由 $2k\le n$ 得到。故所有 $2\le k\le n/2$ 的维数均成立。最小次数 $n=4$ 的允许参数只有 $k=2$，由二维基例覆盖。

若 $2k>n$，令 $k'=n-k$。由 $k+2\le n$ 得 $k'\ge2$，而 $2k'<n$。先对 $k'$ 应用低半结果，再用互补维数对称性得到 $k$ 维零和空间。$n=4$ 时条件只允许 $k=2$，恰由前述二维基例处理。
\end{proof}

\section{奇数次数的维数拼接与参数范围}

除了高维结果，低维基例可产生 $k=6,7,8,9$。若 $E,F$ 分别是 $k$、$\ell$ 维零和子空间，且 $k\ell<n$，形式化的线性代数引理保证存在非零标量 $a$，使 $aF\subseteq P_E(K)$。令
\[
 G=P_E^{-1}(aF).
\]
因为 $\ker P_E=E$、$\dimf P_E(K)=n-k$，有 $\dimf G=k+\ell$。子空间多项式复合公式和系数判据表明 $G$ 仍为零和子空间。因此
\begin{equation}\label{eq:add}
 k,\ell\text{ 已有零和子空间且 }k\ell<n
 \quad\Longrightarrow\quad k+\ell\text{ 也有零和子空间}.
\end{equation}
由 $3,4$ 维基例，依次用 $(3,3),(3,4),(4,4),(3,6)$ 得到 $6,7,8,9$ 维。$k\ge10$ 时用能量结果。这样，对于 $2k+1\le n$ 的低半范围，除 $k=3,4,5$ 外不再需要新基例。

若 $2k+1>n$，令 $k'=n-k$。因 $n$ 奇数且 $k+3\le n$，可验证 $2k'+1\le n$，并且 $k'\ge3$。先在较小的 $k'$ 维数应用低半结论，再用引理\,\ref{lem:complement}取得 $k$ 维零和子空间。该步骤覆盖了奇数次数低半与高半之间的全部参数。

\begin{proof}[定理\,\ref{thm:main} 的证明]
根据 $2k+1\le n$ 或 $2k+1>n$ 分情况。前者中，若 $k=3,4,5$，分别应用命题\,\ref{prop:k3}、定理\,\ref{thm:k4}、定理\,\ref{thm:k5}；若 $6\le k\le9$，应用\eqref{eq:add}；若 $k\ge10$，应用第\,\ref{lem:energy}节的高维结论。后三类结论的阈值均满足所需的 $n\ge2k+1$。后者对 $k'=n-k$ 使用前述低半结论和互补维数对称性。所有构造及拼接均在任意有限域 $K$ 上进行，只使用 $|K|=2^n$ 与特征 $2$，所以保留了定理的全部量词。
\end{proof}

\begin{theorem}[完整奇偶次数形式]\label{thm:full}
设 $n,k\in\mathbb N$，$|K|=2^n$，且 $K$ 为特征 $2$ 的有限域。若
\[
 (\operatorname{Odd}(n)\land 3\le k\land k+3\le n)
 \quad\text{或}\quad
 (\operatorname{Even}(n)\land 2\le k\land k+2\le n),
\]
则 $K$ 含有 $k$ 维零和子空间。
\end{theorem}
\begin{proof}
奇数情形由定理\,\ref{thm:main}给出，偶数情形由定理\,\ref{thm:even}给出。两个结论都对所有满足基数条件的有限域 $K$ 成立。
\end{proof}

\section{Lean 形式化、复现与范围说明}\label{sec:formal}
\sloppy

项目使用 Lean 4 与 mathlib\cite{Lean4,Mathlib}。原始奇数次数命题原样定义为
\[
\forall n,k,\quad \operatorname{Odd}(n)\to 3\le k\to k+3\le n\to
\forall K,\quad |K|=2^n\to \operatorname{HasZeroSumSubspace}(K,k).
\]
{\raggedright
原命题定义于 \texttt{Statement.lean}，奇数次数结论由 \texttt{oddCarletConjecture\_main} 在 \texttt{MainTheorem.lean} 中证明。
偶数定理 \texttt{evenDegreeCarlet\_unconditional} 位于 \texttt{EvenDegree.lean}，适用于偶数 $n$、任意 $|K|=2^n$ 的有限域，以及 $2\le k$、$k+2\le n$。
两项定理合起来给出定理\,\ref{thm:full}。相关证明入口如下：
\par}
\begin{center}
\small\begin{tabular}{>{\raggedright\arraybackslash}p{0.43\textwidth}>{\raggedright\arraybackslash}p{0.47\textwidth}}
\toprule
数学环节 & Lean 文件\\
\midrule
逆元和—根积系数判据、子空间多项式 & \texttt{Elementary.lean}, \texttt{Statement.lean}\\
三维 Kloosterman/Hasse 基例 & \texttt{KloostermanThree.lean}\\
四维奇数次数基例 & \texttt{K4AllOddDegreesProbe.lean} 及其导入模块\\
五维范数曲线构造 & \texttt{K5NormOneConstructionProbe.lean} 及其导入模块\\
概率递推与两步能量收缩 & \texttt{ProbabilityEnergy.lean}, \texttt{TwoStepEnergy.lean}\\
维数拼接、互补维数约化 & \texttt{AdditiveClosure.lean}, \texttt{LowDimensionReduction.lean}\\
偶数二维基例、三维基例与递推 & \texttt{EvenDegree.lean}, \texttt{KloostermanThree.lean}\\
奇数主定理与偶数完整定理 & \texttt{MainTheorem.lean}, \texttt{EvenDegree.lean}\\
\bottomrule
\end{tabular}
\end{center}

{\raggedright
2026 年 11 月 11 日的验收使用 Lean 4.33.0。完整默认构建检查 274 个模块，独立验证审计 942 项声明，并检查 188 个 Hasse 依赖模块。Lean 编译警告为零，源码与 OLean 哈希全部匹配。
奇数主定理与偶数完整定理均通过独立陈述核对。两者的公理集合仅含 \texttt{propext}、\texttt{Classical.choice} 和 \texttt{Quot.sound}；审计不含 \texttt{sorryAx} 或自建未证明公理。
Hasse 界来自固定 AINTLIB 快照\cite{AINTLIB}（\texttt{1c1c7466\allowbreak 4e40071c\allowbreak 2c2165bc\allowbreak 55ca2616\allowbreak a67ccd6b}），源文件为 \texttt{HasseBound.lean}，许可为 Apache 2.0。
复现命令为：
\par}
\begin{verbatim}
cd formalization
python -I build.py
python -I verify_build.py
python -I audit_final_theorem.py
\end{verbatim}
完整验收摘要和定理公理记录分别保存在以下文件中：\par
\begin{quote}\small\texttt{verified\_build\_summary.json}\\
\texttt{final\_theorem\_audit.json}\end{quote}
依赖的 mathlib 源码基线为：\par
\begin{quote}\small\ttfamily db584cd6d46c92f209a44c0f1c829460d327499d\end{quote}

本文的奇数次数论证遵循当前通过 Lean 检查的范数曲线与能量路线；偶数次数论证使用二维、三维基例、维数相加和互补维数对称性。有限计算只用于文中明确限定的有限小次数分支，并已在 Lean 中验证相应构造及其向任意同基数有限域的传递。


\section*{致谢与 AI 使用说明}
写作和代码整理过程中使用了 ChatGPT 辅助。形式化检查由 Lean 内核及项目构建脚本执行；作者对数学内容、引用、代码来源和最终稿件负责。

\sloppy\begin{thebibliography}{9}\small
\bibitem{Carlet2025}
C.~Carlet, \emph{On the vector subspaces of $\mathbb F_{2^n}$ over which the multiplicative inverse function sums to zero}, Designs, Codes and Cryptography 93 (2025), 1237--1254. DOI: \href{https://doi.org/10.1007/s10623-024-01531-6}{10.1007/s10623-024-01531-6}.

\bibitem{Cheng2026}
K.~Cheng, \emph{Counting zero-sum subspaces for the multiplicative inverse function}, arXiv:2610.06445 (2026), \url{https://arxiv.org/abs/2610.06445}.
\bibitem{Silverman}
J.~H.~Silverman, \emph{The Arithmetic of Elliptic Curves}, 2nd ed., Graduate Texts in Mathematics 106, Springer, 2009.

\bibitem{Lean4}
L.~de Moura and S.~Ullrich, \emph{The Lean 4 Theorem Prover and Programming Language}, in: Automated Deduction -- CADE 28, Lecture Notes in Computer Science 12699, Springer, 2021, pp.~625--635.

\bibitem{Mathlib}
The mathlib Community, \emph{The Lean Mathematical Library}, Proceedings of the 9th ACM SIGPLAN International Conference on Certified Programs and Proofs (CPP 2020), pp.~367--381. DOI: \href{https://doi.org/10.1145/3372885.3373824}{10.1145/3372885.3373824}.

\bibitem{AINTLIB}
C.~Birkbeck, \emph{AINTLIB}, Hasse--Weil bound module, project snapshot \texttt{1c1c7466\allowbreak 4e40071c\allowbreak 2c2165bc\allowbreak 55ca2616\allowbreak a67ccd6b} (Apache License 2.0).
\end{thebibliography}
\end{document}
